\section{总结与延伸}

\subsection{本章小结}

\subsection{从四条线收束成一条主线}

这一讲虽然看起来覆盖了很多东西，但其实只有一条核心主线：\textbf{从噪声到样本的变换，最终都可以看成一条可学习的运输路径}。GAN 用对抗博弈逼近这条路径，diffusion / rectified flow 用逐步去噪或速度场积分来实现它，latent diffusion 把这条路径搬到更便宜的表示空间里，Transformer 则负责把条件和序列建模接起来。

\begin{figure}[H]
\centering
\includegraphics[width=0.95\textwidth]{slides-images/slide-121.jpg}
\caption{总图：GAN、diffusion、latent diffusion 和现代条件生成管线已经合流成一套工程体系\protect\footnotemark}
\end{figure}
\footnotetext{Slides 第122页。}

\begin{importantbox}{最后的 takeaways}
\begin{enumerate}
    \item \textbf{GAN} 的核心价值是高质量样本，但它训练不稳，且不适合把自己当作标准监督学习来理解。
    \item \textbf{Diffusion / rectified flow} 把生成拆成多个小步，训练更稳，数学和工程解释都更清楚。
    \item \textbf{CFG} 几乎是现代条件扩散的标配，它让“听话程度”成为可调参数。
    \item \textbf{Latent diffusion} 让高分辨率生成变得可承受，是今天最常见的图像和视频生成范式。
    \item \textbf{Generalized diffusion} 告诉我们：很多名字不同的方法，本质上只是同一套运输问题的不同参数化。
    \item \textbf{AR on discrete latents} 说明 autoregressive 并没有过时，只是从 raw pixels 迁移到了更合适的表示空间。
\end{enumerate}
\end{importantbox}

\subsection{和前两讲的联系}

如果把上一讲和这一讲合起来看，就会发现生成式建模其实一直在做同一件事：用可计算的方式把简单分布变成复杂数据分布。区别只在于，你是选择直接分解概率、通过 latent 变分近似、用对抗博弈、还是用逐步去噪。现代系统之所以复杂，是因为它们已经不满足于“能生成”，而要同时满足“质量高、可控、可扩展、可部署”。

\begin{figure}[H]
\centering
\includegraphics[width=0.95\textwidth]{slides-images/slide-122.jpg}
\caption{下一讲预告：Vision + Language\protect\footnotemark}
\end{figure}
\footnotetext{Slides 第123页。}

\subsection{延伸阅读}

\begin{itemize}
    \item Goodfellow et al., \emph{Generative Adversarial Nets} (2014)
    \item Ho et al., \emph{Denoising Diffusion Probabilistic Models} (2020)
    \item Song et al., \emph{Score-Based Generative Modeling through SDEs} (2021)
    \item Liu et al., \emph{Flow Straight and Fast: Learning to Generate and Transfer Data with Rectified Flow} (2022)
    \item Rombach et al., \emph{High-Resolution Image Synthesis with Latent Diffusion Models} (2022)
    \item Peebles and Xie, \emph{Scalable Diffusion Models with Transformers} (2023)
\end{itemize}

\end{document}
